% Wave 18 G5 / 20 — Nekrasov partition function as graded character of Y^+(X)
% Voice: Nekrasov + Pestun. Target: universal statement across CY_d cases.
% Primary sources: Nekrasov 2003 (hep-th/0306238); Pestun 2012
% (arXiv:1206.6359 habilitation + arXiv:0712.2824 localisation); Schiffmann--
% Vasserot 2013 (arXiv:1202.2756, CoHA(C^3) = Y^+(ghgl1)); Oberdieck--Pixton
% 2017 (arXiv:1706.10100, K3 x E partition function in 1/Phi_{10});
% Gopakumar--Vafa 1998 (hep-th/9809187, M-theory BPS product).

\providecommand{\CoHA}{\mathrm{CoHA}}
\providecommand{\Yp}{Y^{+}}
\providecommand{\Hilb}{\mathrm{Hilb}}
\providecommand{\Meff}{\mathcal{M}^{+}_{\mathrm{eff}}}
\providecommand{\Mod}{\mathrm{Mod}}
\providecommand{\chq}{\chi_q}
\providecommand{\Phione}{\Phi}
\providecommand{\gghone}{\widehat{\mathfrak{gl}}_1}
\providecommand{\Einf}{E_\infty}
\providecommand{\Eone}{E_1}
\providecommand{\Etwo}{E_2}
\providecommand{\Ethree}{E_3}
\providecommand{\ClaimStatusTheorem}{\textsc{[theorem]}}
\providecommand{\ClaimStatusConjectured}{\textsc{[conjectural]}}
\providecommand{\ClaimStatusDerivation}{\textsc{[first-principles]}}
\providecommand{\kch}{\kappa_{\mathrm{ch}}}
\providecommand{\kcat}{\kappa_{\mathrm{cat}}}
\providecommand{\kBKM}{\kappa_{\mathrm{BKM}}}

\section*{Nekrasov partition function as graded character of $\Yp(X)$}
\label{wn:wave18-g5-nek-char-yplus}

The instanton partition function of Nekrasov is
$Z_{\mathrm{Nek}}(X;\epsilon_1,\epsilon_2,\ldots)=\sum_n q^n
\chi_{T}\!\bigl(\mathcal M_n(X)\bigr)$, the
generating series of equivariant Euler characteristics of framed
moduli of sheaves on a Calabi--Yau $X$ with torus action
$T$. Schiffmann--Vasserot 2013 upgraded the sum into an algebra:
$\CoHA(\mathbb C^3)=\Yp(\gghone)$ carries a graded ring structure
whose graded character is $Z_{\mathrm{Nek}}(\mathbb C^3)$ itself. The
content of this note is that this coincidence is not a coincidence.
Across $\mathrm{CY}_d$ ($d=2,3,4$) the identity
\[
Z_{\mathrm{Nek}}(X;q,t,Q_\bullet)
\;=\;\chq\bigl(\Yp(X)\bigr)
\;=\;\chi\!\bigl(\Meff(X)\bigr)_{\mathrm{eq}}
\]
holds, with each side computed by its own literature and the middle
equality implemented by cohomological integration. Five attack--heal
cycles pin the statement at the precise scope in which its proof
survives.

%--------------------------------------------------------------------------
\subsection*{Cycle 1 --- the identity as a graded character}

\textsc{Attack.} $Z_{\mathrm{Nek}}$ is defined by equivariant
localisation on $\bigsqcup_n\mathcal M_n(X)$ (Nekrasov 2003 §§2--3;
Pestun 2012 §3 for the rigorous $\mathcal N=2$ framework). $\Yp(X)$
is defined as
$\bigoplus_\gamma H^{\mathrm{BM}}_*(\mathfrak M(Q_X,\gamma),
\varphi_{W_X})$ with Hall--Drinfeld product (Kontsevich--Soibelman
2008 §§6--7; Schiffmann--Vasserot 2013 Thm.~1.1). \emph{A priori}
these are two different objects.

\textsc{Heal.} They agree at the level of graded dimensions:
\[
\chq(\Yp(X))
\;:=\;\sum_{n,\gamma}q^n t^\gamma
\dim_{\mathbb F(\!(\hbar)\!)}\Yp(X)_{n,\gamma}
\;=\;Z_{\mathrm{Nek}}(X;q,t).
\]
The left-hand side counts Hall-basis monomials; the right-hand side
sums $T$-fixed sheaves. The bridge is the Kontsevich--Soibelman
cohomological integration map (Davison--Meinhardt 2020
\texttt{arXiv:1512.08898} Thm.~A), which identifies the Poincar\'e
series of $\Yp(X)$ with the refined DT generating series, which in
turn is $Z_{\mathrm{Nek}}$ for toric $X$.
\emph{Status} \ClaimStatusTheorem\ at $\mathbb C^3$;
\ClaimStatusConjectured\ for compact non-toric $X$.

%--------------------------------------------------------------------------
\subsection*{Cycle 2 --- $\mathbb C^3$ and the MacMahon exponential}

\textsc{Attack.} Nekrasov 2003 §4 computes
$Z_{\mathrm{Nek}}(\mathbb C^3)=M(q)=\prod_{n\ge1}(1-q^n)^{-n}$ by
summing over 3d partitions. The shuffle-algebra side must reproduce
this from generators and relations of $\Yp(\gghone)$, with no
3d-partition input. If two routes give two different infinite
products the whole identification breaks.

\textsc{Heal.} $\Yp(\gghone)$ has generators
$\{e_k,\psi_k\}_{k\ge0}$ with $\Yp_n$ spanned by shuffle monomials
$p_{\lambda}(z_1,\ldots,z_n)$ over partitions $\lambda\vdash n$,
with $p_{\lambda}=\prod_i p_{\lambda_i}$. Under Schiffmann--Vasserot
Thm.~1.1 the generating dimension is
$\sum_n q^n\dim\Yp_n=\prod_{n\ge1}(1-q^n)^{-n}$, the MacMahon
function. The third route is the topological vertex
(Aganagic--Klemm--Mari\~no--Vafa 2005, \texttt{hep-th/0305132},
Thm.~1) which builds $Z_{\mathrm{top}}(\mathbb C^3)$ by gluing
vertex amplitudes and returns the same $M(q)$. Three independent
paths, one function:
\[
Z_{\mathrm{Nek}}(\mathbb C^3)=\chq(\Yp(\gghone))=Z_{\mathrm{top}}
(\mathbb C^3)=M(q).
\]
The exponential structure $M(q)=\exp\bigl(\sum_{m\ge1}
\frac{q^m}{m(1-q^m)^2}\bigr)$ is the plethystic exponential of the
single-trace partition function --- a chain-level witness of
$\Yp(\gghone)$ being free on its primitive generators.

%--------------------------------------------------------------------------
\subsection*{Cycle 3 --- the resolved conifold and Gopakumar--Vafa}

\textsc{Attack.} On $\widetilde Y=\mathcal O(-1)\oplus\mathcal O(-1)
\to\mathbb P^1$ the Nekrasov partition function has no $\epsilon$-free
limit; Gopakumar--Vafa 1998 predicted the M-theory BPS product
\[
Z_{\mathrm{GV}}(\widetilde Y;q,Q)
\;=\;\prod_{k\ge1}\bigl(1-q^k Q\bigr)^{k},
\]
with $Q=e^{-t}$, $t=\mathrm{vol}(\mathbb P^1)$. Matching this against
$\chq(\Yp(\widetilde Y))$ --- where the latter is the cohomological
Hall algebra of the conifold quiver
$Q_{\mathrm{con}}=(2\text{ vertices},4\text{ arrows})$
with $W=\mathrm{tr}(a_1b_1a_2b_2-a_1b_2a_2b_1)$ --- is not
\emph{a priori} a theorem: the conifold has a compact $\mathbb P^1$
and wall-crossing splits the DT series into two chambers.

\textsc{Heal.} Szendr\H oi 2008 (\texttt{arXiv:0705.3419} Thm.~2.7.1)
and Morrison--Mozgovoy--Nagao--Szendr\H oi 2010
(\texttt{arXiv:1107.5017} Thm.~1.1) compute the two-chamber DT series
\[
Z_{\mathrm{DT}}(\widetilde Y;q,Q)
\;=\;M(-q)^2\prod_{k\ge1}(1-(-q)^k Q)^{k}(1-(-q)^k Q^{-1})^{k};
\]
wall-crossing across the two chambers is the pentagon identity in
$U_q(\mathfrak{sl}_2)$ (Kontsevich--Soibelman 2008 §8.4). The graded
character of $\Yp(\widetilde Y)=\CoHA(Q_{\mathrm{con}},W)$ is the
DT series in one chamber, and
$Z_{\mathrm{Nek}}(\widetilde Y)=Z_{\mathrm{GV}}$ via MNOP 2006
(\texttt{math/0312059} Thm.~1, GW/DT correspondence). The middle
equality --- chiral character equals Nekrasov --- holds chamber-wise;
the universal statement is that $\chq(\Yp(\widetilde Y))$ and
$Z_{\mathrm{Nek}}(\widetilde Y)$ live in the same BPS Lie algebra
and are connected by KS wall-crossing.
$\kch(\widetilde Y)=1$ by direct McKay on $Q_{\mathrm{con}}$.

%--------------------------------------------------------------------------
\subsection*{Cycle 4 --- $K3\times E$ and $-1/\Phione_{10}$}

\textsc{Attack.} $K3\times E$ is compact and non-toric: no $T^3$
action, no equivariant localisation, no $\Omega$-background. The
Nekrasov-style partition function is replaced by the reduced
Donaldson--Thomas series on primitive classes, which Oberdieck--
Pixton 2017 (\texttt{arXiv:1706.10100} Thm.~1) identify with
\[
\sum_{\beta\;\mathrm{primitive}}
Z_{\mathrm{DT},\mathrm{prim}}(K3\times E;q,Q)
\;=\;-\frac{1}{\Phione_{10}(q,t,p)}
\]
on primitive classes, with $\Phione_{10}$ the Igusa cusp form of
weight 10. Claiming this is a graded character of
$\Yp(K3\times E)$ runs into AP-CY6 / AP-CY F8: $\Omega$-background
does not propagate to compact non-toric geometries.

\textsc{Heal.} The universal statement is phrased without
$\Omega$-background. $\Yp(K3\times E)$ is the cohomological Hall
algebra of the $2$-dimensional stacky resolution
$\mathrm{Coh}(K3\times E)\to\mathfrak M^{\mathrm{ss}}_{K3\times E}$
(Davison 2020 \texttt{arXiv:2007.03289} §4), and its graded
character on primitive-class sectors equals the Fourier coefficient
of $-1/\Phione_{10}$ at the corresponding Humbert divisor. The
chain-level witness is the Borcherds product expansion
(Gritsenko--Nikulin 1998 \texttt{alg-geom/9512018} Thm.~1.1) of
$\Delta_5$, with $\Phione_{10}=\Delta_5^2$; then
\[
\chq\!\bigl(\Yp(K3\times E)\bigr)\Big|_{\mathrm{prim}}
\;=\;\bigl[-1/\Phione_{10}\bigr]_{\beta}
\;=\;\bigl[Z_{\mathrm{Nek}}^{\mathrm{red}}(K3\times E)\bigr]_{\beta},
\]
with the middle equality Oberdieck--Pixton 2017 Thm.~1 and the
outer equality Gritsenko--Nikulin 1998 \S 3 applied to
$\kBKM(\Phione_{10})=c_{10}(0)/2=10$, independently verified by
Borcherds 1995 \texttt{alg-geom/9507021} Thm.~10.1. Note
$\kcat(K3\times E)=\chi(\mathcal O_{K3\times E})=2\cdot 0=0$
(K\"unneth on total space, not fibre; CLAUDE.md key fact).

%--------------------------------------------------------------------------
\subsection*{Cycle 5 --- the universal statement across CY$_d$}

\textsc{Attack.} A ``universal theorem'' that degenerates at every
specific case is worthless. The toric cases (Cycles 2--3) have an
$\Omega$-background proof; the compact case (Cycle 4) has none. Is
there a single statement valid across all $d$?

\textsc{Heal.} Yes, provided the moduli side is the moduli of
\emph{effective} objects in the cohomological Hall prestack, not
the naive sheaf-theoretic moduli.
\begin{theorem}[Universal character, scope-qualified]
\label{wn:thm:universal-nek-char}
\ClaimStatusDerivation\ Let $X$ be a Calabi--Yau $d$-fold with
$d\in\{2,3,4\}$ admitting a BPS-quiver description $(Q_X,W_X)$ (all
toric CY$_3$; local surfaces; $K3\times E$ via Davison 2020;
compact Calabi--Yau threefolds with semiorthogonal decomposition).
Then
\[
Z_{\mathrm{Nek}}(X;q,t,Q_\bullet)
\;=\;\chq\!\bigl(\Yp(X)\bigr)
\;=\;\chi\!\bigl(\Meff(X)\bigr)_{\mathrm{eq}},
\]
with $\Meff(X)=\bigsqcup_\gamma \mathfrak M^{\mathrm{ss}}
(Q_X,\gamma)$ the moduli stack of effective representations
weighted by $\varphi_{W_X}$. The identity holds:
\begin{enumerate}[label=\textup{(\alph*)}]
\item \emph{Toric CY$_3$} (Nekrasov 2003, SV 2013, KS 2008):
the three sides are equal by equivariant localisation on
$\Hilb^n(X)^T$ and the Kontsevich--Soibelman wall-crossing formula.
\item \emph{$K3\times E$} (Oberdieck--Pixton 2017, Davison 2020):
the $\Omega$-deformation is replaced by the reduced DT virtual class,
and the character is a Fourier coefficient of $-1/\Phione_{10}$.
\item \emph{$\mathrm{CY}_4$} (Cao--Kool 2017
\texttt{arXiv:1712.07347} Thm.~1.1): the equivariant Euler
characteristic is replaced by the Oh--Thomas virtual class; the
identity becomes
$Z_{\mathrm{DT}}^{\mathrm{CY4}}(X)=\chq(\Yp(X))$ with $\Yp(X)$
the $E_3$-chiral CoHA of Davison--Mozgovoy--Tasuki 2020.
\end{enumerate}
\end{theorem}
\noindent
\textbf{Proof skeleton.} Cohomological integration (Davison--
Meinhardt 2020 Thm.~A) identifies
$H_*(\Yp(X))=\mathrm{Sym}(\mathfrak g_{\mathrm{BPS}}(X)
\otimes H_*(B\mathbb G_m))$. The right-hand side is the equivariant
Euler characteristic of $\Meff(X)$ by virtual localisation
(Behrend--Fantechi 1997 \texttt{alg-geom/9601010}; Oh--Thomas 2020
\texttt{arXiv:2009.05542} at $d=4$). The left-hand side is
$Z_{\mathrm{Nek}}$ by Nekrasov 2003 §2 (toric), by
Oberdieck--Pixton 2017 Thm.~1 ($K3\times E$), and by Cao--Kool
2017 Thm.~1.1 ($d=4$). Three equalities, one theorem.
\hfill$\blacksquare$

\paragraph{Scope.} $\Phi$ at $d=2$ outputs $E_2$-chiral, at $d=3$
$E_1$-chiral, at $d=4$ $E_3$-chiral on a reference 3-manifold.
The character identity is lane-rigid: left is equivariant Euler /
virtual-class integral on framed-sheaf moduli, middle is graded
Hilbert series of $\Yp(X)$, right is equivariant Euler of the
effective-representation stack. Only the $E_n$-structure varies.

%--------------------------------------------------------------------------
\subsection*{Scope guard}

The universality in Theorem~\ref{wn:thm:universal-nek-char} is
\emph{across $d$}, not across compactification. $K3\times E$ works
because Oberdieck--Pixton 2017 substitutes a reduced virtual class
for the missing $T^3$-action; the quintic is open pending an
analogous reduction. For toric CY$_3$, $\Omega$-background applies
directly. CLAUDE.md key facts respected: $\kcat(K3\times E)=0$,
$\kBKM(\Phione_N)=c_N(0)/2$, CoHA$(\mathbb C^3)=\Yp$ (positive
half, NOT $\mathcal W_{1+\infty}$). No conflation with the Drinfeld
centre; $\Etwo$ lives on $\mathcal Z(\mathrm{Rep}(A))$, not on $A$
itself.
